\documentclass[10pt,a4paper]{amsart}
\usepackage[margin=19mm]{geometry}
\usepackage{amsmath,amssymb,mathtools}
\usepackage[scr=rsfs,cal=euler]{mathalfa}
\usepackage{xfrac}
\usepackage[unicode,hidelinks,bookmarks=false]{hyperref}
\newcommand{\ie}{i.e.\ }
\newcommand{\cf}{cf.\ }
\newcommand{\resp}{respectively\ }
\newcommand{\Subsection}{\subsection}
\providecommand{\nobreakdash}{\nobreak\hbox{-}}
\setlength{\emergencystretch}{2em}
\multlinegap0pt
\usepackage{amssymb}
\usepackage[shortlabels]{enumitem}
\newtheorem{lemma}{Lemma}
\newtheorem{definition}{Definition}
\newtheorem{theorem}{Theorem}
\renewcommand{\epsilon}{\varepsilon}
\newcommand{\fm}{\mathfrak{m}}
\renewcommand{\leq}{\leqslant}
\title{Asymptotics of high-codimensional area-minimizing currents in hyperbolic space}
\author{Xumin Jiang, Jiongduo Xie}
\date{Source: arXiv:2601.04027v3 (CC BY 4.0)}
\begin{document}
\maketitle

\noindent\emph{Source and licence.} Asymptotics of high-codimensional area-minimizing currents in hyperbolic space. Version arXiv:2601.04027v3 (CC BY 4.0), distributed by arXiv under the Creative Commons Attribution 4.0 International licence, http://creativecommons.org/licenses/by/4.0/, as recorded in the arXiv licence metadata for this exact version and shown on the abstract page https://arxiv.org/abs/2601.04027v3. Full licence notice: \texttt{licenses/arxiv-2601.04027-CC-BY-4.0.txt}.\par\smallskip

\noindent\textit{Editorial scope.} Counted unit: the article's theorem thm-c41 (source0.tex lines 1206--1219). The article's complete original proof of it is delivered with it (source0.tex lines 1221--1352, one \texttt{proof} environment of 132 lines). No mathematics is written, repaired or re-derived: the statement and the proof are the article's own lines, in the article's own notation, and no retained line is substituted or re-flowed.  Retained uncounted as the source-local context the proof needs: lines 137--139; lines 243--245; lines 786--788; lines 1018--1033; lines 1082--1096; lines 1138--1149; lines 1152--1154; lines 1164--1166; lines 1172--1174.  Nothing was omitted from any retained span, so the package carries no omission record.  No retained line contains a citation, so the wrapper carries no bibliography and no bibliography entry.  No retained line contains a figure environment, an image-inclusion command or drawing code, so no image is delivered and the package declares no asset dependency.  Editorial formatting only, in addition: an AMS \texttt{amsart} preamble that restates the article's own declarations and macros used by the retained lines, namely the environments \texttt{lemma}, \texttt{definition}, \texttt{theorem} on separate counters, together with the standard packages those lines need.  The retained environments print in this document as Lemma 1, Definition 1, Theorem 1 in order of appearance, whereas the article prints the same items with the numbering of its own full document.  The complete native source of the article as distributed in the arXiv e-print for this exact version is bundled unchanged as \texttt{sources/192257/source0.tex}.
\begin{align}\label{eq-H+-0}
    H^+ = \bigl\{(x', x^\fm)\in \mathbb{R}^\fm: (x',0) \in T_Q\Gamma,\ x^\fm \in \mathbb{R}^+\bigr\}.
\end{align}

\begin{align}\label{eq-u-exp-0}
    \mathbf{u} = \varphi(y') +\sum_{i=2}^{n}\mathbf{c}_i(y')(y^n)^i +\sum_{i=n+1}^k\sum_{j=0}^{\lfloor \frac{i-1}{n}\rfloor} \mathbf{c}_{i,j}(y')(y^n)^i(\log(y^n))^j +\mathbf{R}_k,
\end{align}

\begin{align}\label{eq-main}
    g^{ij}\frac{\partial^2 u_s}{\partial y^i \partial y^j} - \frac{n}{y^{n}} \frac{\partial u_s}{\partial y^{n}} =0,
\end{align}

\begin{lemma}[Formal Series Solution]\label{lem-Form-Comp}
    Let $\varphi \in C^\infty(B_r')$ be a smooth function. Then there exist smooth $(\fm-n)$-valued functions $\mathbf{c}_i$ and $\mathbf{c}_{i,j}$ on $B_r'$ (with $(i,j)\neq (n+1,0)$) that depend on $\varphi$ and the first global coefficient $\mathbf{c}_{n+1,0}$ such that, for any given smooth $(\fm-n)$-valued function $\mathbf{c}_{n+1,0}$ on $B_r'$, there exists a unique formal series of the form
    \begin{align}\label{eq-expan}
        \mathbf{u} = \varphi +\sum_{i=2}^{n}\mathbf{c}_i(y')(y^n)^i +\sum_{i=n+1}^\infty \sum_{j=0}^{\left\lfloor \frac{i-1}{n}\right\rfloor} \mathbf{c}_{i,j}(y')(y^n)^i\bigl(\log (y^n)\bigr)^j \quad \text{in } B^+_r,
    \end{align}
    which formally solves the minimal surface system \eqref{eq-main}. Here  $\mathbf{c}_i(y')=\mathbf{0}$ whenever $i\leq n$ is an odd integer.

    In addition, we define the partial sum $\mathbf{T}_k$ by
    \begin{align}\label{eq-uk}
        \mathbf{T}_k = \varphi +\sum_{\substack{i=2 \\ i\ \text{even}}}^{2\left\lfloor \frac{n}{2}\right\rfloor}\mathbf{c}_i(y') (y^n)^i +\sum_{i=n+1}^k\sum_{j=0}^{\left\lfloor \frac{i-1}{n}\right\rfloor} \mathbf{c}_{i,j}(y') (y^n)^i\bigl(\log (y^n)\bigr)^j.
    \end{align}
    This partial sum satisfies the asymptotic estimate
    \begin{align}
        \mathcal{Q}(\mathbf{T}_k)=O\Bigl((y^n)^{k}\bigl(\log (y^n)\bigr)^{\left\lfloor \frac{k-1}{n}\right\rfloor}\Bigr).
    \end{align}
\end{lemma}

\begin{definition}\label{def-exp}
A function $\mathbf{w}$ (which may be single-valued or $(\fm-n)$-valued) is said to admit an expansion of order $t^k$ in $B^+_r$ if there exist smooth functions $\mathbf{c}_{i,j}(y')$ on $B_r'$, positive integers $N_i\in \mathbb{N}$, and a constant $\epsilon>0$ such that the decomposition
\begin{align}\label{eq-Taylor}
    \mathbf{w} = \mathbf{T}_k +\mathbf{R}_k := \sum_{i=0}^k \sum_{j=0}^{N_i} \mathbf{c}_{i,j} t^i (\log t)^j+\mathbf{R}_k
\end{align}
holds in $\overline{B^+_r}$. Here, the remainder term $\mathbf{R}_k$ satisfies two key conditions:
\begin{align}\label{eq-Rk-Cke}
 \mathbf{R}_k \in O(t^{k+\epsilon})\cap C^{k,\epsilon}(\overline{B^+_r}),
\end{align}
and for all nonnegative integers $p, q \in \mathbb{N}$, the tangential derivative estimate
\begin{align}\label{eq-Rk-Tan}
    \|D_{y'}^p \mathbf{R}_k \|_{C^{q,\alpha}_h(G_r)}\leq C_{p,q,k}\,t^{k+\epsilon}
\end{align}
is valid, where $C_{p,q,k}$ denotes a constant independent of $t$ and $y'$.
\end{definition}

\begin{lemma}\label{lem-c2} Assume the same hypotheses as in Lemma \ref{lem-Form-Comp}. Then the $s$ component of $\mathbf c_{2}$ is
\begin{align}\label{eq-c-s2}
    c_{s,2}(y') = \frac{1}{2}\sqrt{1+|D\varphi_s(y')|^2} (H(y'), \nu_s(y')),
\end{align}
where $H$ is the mean curvature vector of $\Gamma$ at 
\begin{align}
 (y', 0, \varphi(y')) = (y^1,\cdots, y^{n-1}, 0, \varphi((y^1,\cdots, y^{n-1})),  
\end{align} and $\nu_s$ is a unit normal vector in the form of
\begin{align}
    \nu_s = \frac{(-D\varphi_s, \mathbf e_s)}{\sqrt{1+|D \varphi_s|^2}}.
\end{align}
\end{lemma}

    \begin{align}\label{eq-main-comp1}
        u_{s, tt} -\frac{n}{t} u_{s,t} = \frac{n}{t}((g^{nn})^{-1}-1)u_{s, t} -(g^{nn})^{-1}\sum_{(i,j)\neq (n,n)}g^{ij}u_{s,ij},
    \end{align}

\begin{align}\label{eq-c2-local}
    c_{s,2} =\frac{1}{2n-2}\sum_{1\leq \alpha,\beta<n}g^{\alpha\beta} \varphi_{s,\alpha \beta}. 
\end{align}

\begin{align}\label{eq-II}
\mathrm{II}_{\alpha \beta} &= \left(\nabla_{\mathbf{F}_{,\alpha}}\mathbf{F}_{,\beta}\right)^\perp = \left(\sum_{\xi=1}^{\fm-n} \varphi_{\xi, \alpha\beta}\partial_{\xi}\right)^\perp.
\end{align}

\begin{theorem}\label{thm-c41}
Assume the same hypotheses as in Lemma \ref{lem-Form-Comp}. If in addition $n=3$, then the $s$-component of $\mathbf{c}_{4,1}$ is
\begin{align}\label{eq-c-41}
c_{s, 4,1} = -\frac{1}{8}\sqrt{1+|D\varphi_s(y')|^2} \left( (\Delta_\Gamma H)^\perp-2|H|^2 H + 2Q(\operatorname{II})(H), \nu_s \right),
\end{align}
where $H$ and $\operatorname{II}$ denote the mean curvature vector and the second fundamental form of $\Gamma$ at
\begin{align}
(y', 0, \varphi(y')) = \bigl(y^1,\cdots, y^{n-1}, 0, \varphi(y^1,\cdots, y^{n-1})\bigr),
\end{align}
and $\nu_s$ is the unit normal vector given by
\begin{align}
\nu_s = \frac{(-D\varphi_s, \mathbf{e}_s)}{\sqrt{1+|D \varphi_s|^2}}.
\end{align}
\end{theorem}

\begin{proof}
    First, we rewrite \(c_{s,2}\). By \eqref{eq-II},
\begin{align}\label{eq-H-exp}
    (n-1)H = \left(g^{\alpha\beta}\sum_{\xi=1}^{\fm-n} \varphi_{\xi, \alpha\beta}\partial_{\xi}\right)^\perp
    = g^{\alpha\beta}\sum_{\xi=1}^{\fm-n} \varphi_{\xi, \alpha\beta}\partial_{\xi}
    - g^{\alpha\beta} (\varphi_{\alpha\beta}\cdot \varphi_\gamma)g^{\gamma\eta}\left(\partial_\eta+ \sum_{\xi=1}^{\fm-n} \varphi_{\xi, \eta} \partial_{\xi}\right),
\end{align}
where \(\alpha, \beta, \gamma, \eta\) are summed from \(1\) to \(n-1\). By \eqref{eq-c2-local} and \eqref{eq-H-exp},
\begin{align}\label{eq-c2-H}
    c_{s,2} = \frac{1}{2}(H, \partial_s) +\frac{1}{2n-2} g^{\alpha\beta} (\varphi_{\alpha\beta}\cdot \varphi_\gamma)g^{\gamma\eta}\varphi_{s, \eta}.
\end{align}

Next, we compute \(c_{s, 4,1}\) under a special coordinate system. Let \(Q\in \Gamma\) be the origin of the coordinate system \(y\), which satisfies
\begin{align}
    y^{n+1} = \cdots = y^\fm = 0
\end{align}
on the vertical tangent plane \(H^+\) associated with \(Q\). Clearly, the corresponding defining function \(\varphi\) of \(\Gamma\) satisfies, for \(s=1,\dots, \fm-n\),
\begin{align}\label{eq-Dvarphi-0}
    D\varphi_s (0) = 0.
\end{align}

At the point where
\begin{align}
    y' = (y^1,\dots, y^{n-1}) = 0,
\end{align}
substitute \(u = \varphi + \mathbf{c}_2(y') t^2\) into the right-hand side of \eqref{eq-main-comp1} to obtain the \(O(t^2)\) terms as the coefficient of \(t^2\) in the expression
\begin{align}
    n\left(4\sum_{\xi=1}^{\fm-n} c_{\xi, 2}^2t^2\right)2c_{s,2}
    - \left(1+4 \sum_{\xi=1}^{\fm-n} c_{\xi,2}^2 t^2\right)\Delta' \left(\varphi_s +c_{s,2}t^2\right),
\end{align}
where at \(Q\),
\begin{align}
    \Delta' c_{s,2} = \frac{1}{2}(\Delta'H,\partial_s) + (H_{\xi} \cdot \operatorname{II}_{\xi,\gamma\eta}) \operatorname{II}_{s,\gamma\eta}
\end{align}
by \eqref{eq-c2-H}. Then the \(O(t^2)\) terms on the right-hand side of \eqref{eq-main-comp1} are
\begin{align}\label{eq-c-4-RHS}
    |H|^2 H_s - \frac{1}{2}(\Delta'H,\partial_s) - (H_{\xi} \cdot \operatorname{II}_{\xi,\gamma\eta}) \operatorname{II}_{s,\gamma\eta}.
\end{align}
Thus we obtain at \(Q\) that
\begin{align}\label{eq-c41-Q}
    c_{s,4,1} = -\frac{1}{8} \left( (\Delta'H,\partial_s)-2|H|^2 H_s + 2(H_{\xi} \cdot \operatorname{II}_{\xi,\gamma\eta}) \operatorname{II}_{s,\gamma\eta}\right).
\end{align}

Thirdly, we consider the general case. By \eqref{eq-u-exp-0}, under the special coordinate system $y$ satisfying \eqref{eq-Dvarphi-0},
\begin{align}\label{eq-uQ-y}
    \mathbf u_Q(y', t) = \sum_{i=0}^k\sum_{j=0}^{\lfloor \frac{i-1}{n}\rfloor} \mathbf{d}_{i,j}(y')\, t^i (\log t)^j + \tilde{\mathbf R}_k,
\end{align}
in the sense of Definition \ref{def-exp}. Clearly, $\varphi = \mathbf d_{0,0}$, and
\begin{align}\label{eq-uQ-y-2}
    \mathbf u_Q(y', t) = \varphi(y') + \mathbf{d}_{2,0}(y')t^2 + \sum_{i=4}^k\sum_{j=0}^{\lfloor \frac{i-1}{n}\rfloor} \mathbf{d}_{i,j}(y')\, t^i (\log t)^j + \tilde{\mathbf R}_k.
\end{align}

Now fix a point $P\in \Gamma$ near $Q$, and take a special coordinate system $z$ at $P$ such that at $(z', t) = (0, 0)$,
\begin{align}
    D_{z'} \mathbf u_Q = 0.
\end{align}
By \eqref{eq-u-exp-0}, there exist functions $\mathbf c_{i,j}(y')$ such that
\begin{align}\label{eq-uP-exp}
    \mathbf u_P(z', t) = \sum_{i=0}^k\sum_{j=0}^{\lfloor \frac{i-1}{n}\rfloor} \mathbf{c}_{i,j}(z')\, t^i (\log t)^j + \mathbf R_k
\end{align}
holds under the coordinate system $z$. Our goal is to evaluate $\mathbf c_{4,1}$ at $Q$ for $n=3$.

Since $z$ and $y$ differ by a translation and a linear transformation, there exists an $\fm\times \fm$ matrix $B$ such that
\begin{align}
    y = B(z - z(Q)).
\end{align}
In components, for $i=1,\dots,\fm$,
\begin{align}
    y_i = \sum_{l=1}^\fm B_{il}\,(z-z(Q))_l.
\end{align}
We assume that $y$ and $z$ are Euclidean coordinates in $\mathbb{R}^\fm$ with $y^n=z^n=t$, which implies
\begin{align}\label{eq-B-in}
    B_{in} = \delta_{in}.
\end{align}
Using \eqref{eq-B-in} and \eqref{eq-uP-exp}, we obtain on the minimal surface that, for $s>n$,
\begin{align}
y_s &= \sum_{l=1}^\fm B_{sl}\,(z-z(Q))_l \nonumber\\
&= \sum_{l=1}^{n-1} B_{sl}\,(z-z(Q))_l - \sum_{l=1}^\fm B_{sl}\,z_l(Q) + \sum_{\xi=n+1}^\fm B_{s\xi}\,z_\xi \nonumber\\
&= \sum_{l=1}^{n-1} B_{sl}\,(z-z(Q))_l - \sum_{l=1}^\fm B_{sl}\,z_l(Q) \nonumber\\
&\quad + \sum_{\xi=n+1}^\fm B_{s\xi} \left( \sum_{i=0}^k\sum_{j=0}^{\lfloor \frac{i-1}{n}\rfloor} c_{\xi,i,j}(z')\, t^i (\log t)^j + R_{\xi,k} \right). \label{eq-uQ-z}
\end{align}

Fix $z' = z'(Q)$ and vary $t$ to obtain a vertical ray $L^+ \subseteq H_P^+$, defined as in \eqref{eq-H+-0} with $Q$ replaced by $P$. Then the graph map
\begin{equation}
    \bigl(z',\, t,\, \mathbf{u}_P(z', t)\bigr)
\end{equation}
maps $L^+$ onto a curve contained in the minimal surface. We take derivatives with respect to $t$ along $L^+$ and compute
\begin{align}
    \lim_{t\rightarrow 0}\left(\frac{1}{\log t} \frac{\partial^4 y_s}{\partial t^4} \right)
    = \lim_{t\rightarrow 0}\left(\frac{1}{\log t} \frac{\partial^4 u_{Q,s}}{\partial t^4} \right)
\end{align}
in two different ways. Applying \eqref{eq-uQ-z}, we derive
\begin{align}
    \lim_{t\rightarrow 0}\left(\frac{\partial^2 y_s}{\partial t^2} \right)
    &= 2 \sum_{\xi=n+1}^\fm B_{s\xi} \, c_{\xi,2,0}(Q), \label{eq-c2-trans}\\[4pt]
    \lim_{t\rightarrow 0}\left(\frac{1}{\log t} \frac{\partial^4 y_s}{\partial t^4} \right)
    &= 24 \sum_{\xi=n+1}^\fm B_{s\xi} \, c_{\xi,4,1}(Q). \label{eq-c41-trans}
\end{align}

On the other hand, $y_s = u_{Q, s}(y', t) = u_{Q, s}(y'(z), t)$. To apply \eqref{eq-uQ-y}--\eqref{eq-uQ-y-2}, we compute
\begin{align}
    \frac{\partial d_{s, i,j}}{\partial t} &= \sum_{\alpha=1}^{n-1}\sum_{\xi =n+1}^\fm \frac{\partial d_{s, i,j}}{\partial y^\alpha} \frac{\partial y^\alpha}{\partial z^\xi} \frac{\partial z^\xi}{\partial t} = \sum_{\alpha=1}^{n-1}\sum_{\xi =n+1}^\fm \frac{\partial d_{s, i,j}}{\partial y^\alpha} B_{\alpha\xi} \frac{\partial z^\xi}{\partial t}, \\
    \frac{\partial^2 d_{s, i,j}}{\partial t^2} &= \sum_{\alpha,\beta=1}^{n-1}\sum_{\xi,\zeta =n+1}^\fm \frac{\partial^2 d_{s, i,j}}{\partial y^\alpha \partial y^\beta} B_{\alpha\xi} \frac{\partial z^\xi}{\partial t} B_{\beta\zeta} \frac{\partial z^\zeta}{\partial t}+\sum_{\alpha=1}^{n-1}\sum_{\xi =n+1}^\fm \frac{\partial d_{s, i,j}}{\partial y^\alpha} B_{\alpha\xi} \frac{\partial^2 z^\xi}{\partial t^2}.
\end{align}
For $(i,j)=(0,0)$, we have $d_{s,0,0}=\varphi_s$, and at $Q$,
\begin{align}
    \frac{\partial d_{s, 0,0}}{\partial t}=0, \quad \frac{\partial^2 d_{s, 0,0}}{\partial t^2} (y', 0)=0,
\end{align}
since as $t\rightarrow 0$,
\begin{align}
    \frac{d z^\xi}{d t} = 2c_{\xi, 2, 0} t+o(t) \rightarrow 0.
\end{align}
Then at $Q$, by \eqref{eq-uQ-y-2} and \eqref{eq-c2-trans},
\begin{align}
   \lim_{t\rightarrow 0}\left(\frac{\partial^2 y_s}{\partial t^2} \right) = 2 d_{s, 2, 0}(Q)=2 \sum_{\xi=n+1}^\fm B_{s\xi} c_{\xi,2,0}(Q).
\end{align}
Thus the matrix $(B_{s\xi})_{n+1\leq s, \xi\leq \fm}$ maps $\mathbf c_{2,0}(Q)$ to $\mathbf d_{2,0}(Q) = \frac{1}{2}H(Q)$, where $\mathbf c_{2,0}(Q)$ is given by Lemma \ref{lem-c2} with $\varphi$ replaced by $\varphi_P := \mathbf d_{0,0}$ and $y'$ replaced by $z'(Q)$.

Since the fourth-order derivative of $\varphi(y')+ \mathbf d_{2, 0} t^2$ with respect to $t$ along the ray $L^+$ is bounded, we obtain
\begin{align}
    \lim_{t\rightarrow 0}\left(\frac{1}{\log t} \frac{\partial^4 y_s}{\partial t^4} \right) = 24 d_{s,4,1}(Q).
\end{align}
By \eqref{eq-c41-trans}, it follows that
\begin{align}
    d_{s,4,1}(Q) = \sum_{\xi=n+1}^\fm B_{s\xi} c_{\xi,4,1}(Q).
\end{align}
Hence the matrix $(B_{s\xi})_{n+1\leq s, \xi\leq \fm}$ also maps $\mathbf c_{4,1}(Q)$ to
\begin{align}
\mathbf d_{4,1}(Q) =  -\frac{1}{8} \left( \Delta_\Gamma H-2|H|^2 H + 2(H_{\xi} \cdot \operatorname{II}_{\xi,\gamma\eta}) \operatorname{II}_{\cdot,\gamma\eta}\right),
\end{align}
where $H$, $\Delta_\Gamma H$, and $\operatorname{II}$ are evaluated at $Q$. We then apply Lemma \ref{lem-c2}  to complete the proof.
\end{proof}

\end{document}
